\documentclass[aspectratio=169]{beamer}

\usepackage[utf8]{inputenc}
\usepackage[T1]{fontenc}
\usepackage{helvet}
\renewcommand{\familydefault}{\sfdefault}
\usepackage{amsmath, amssymb}
\usepackage{booktabs}
\usepackage{tcolorbox}
\usepackage{array}
\usepackage{xcolor}
\usepackage{hyperref}
\usepackage{tikz}

\tcbset{boxsep=2pt, left=4pt, right=4pt, top=2pt, bottom=2pt}

% Colores UAEM
\definecolor{uaemcolor}{RGB}{0, 80, 100}
\definecolor{grisclaro}{RGB}{245, 245, 245}
\definecolor{verdeenv}{RGB}{34, 120, 70}
\definecolor{naranjaaqi}{RGB}{200, 90, 20}

\usetheme{default}
\usecolortheme{default}
\setbeamertemplate{navigation symbols}{}
\setbeamercolor{normal text}{fg=black}
\setbeamercolor{frametitle}{fg=black}

% Frametitle estilo SimplePlus: barra de color arriba, texto negro, limpio
\setbeamertemplate{frametitle}{
  \nointerlineskip

  \vskip0.3cm
  \hspace{0.5cm}{\large\bfseries\insertframetitle}
  %\vskip0.15cm
  \hspace{0.5cm}\textcolor{uaemcolor!60}{\rule{\dimexpr\textwidth-0.0cm}{0.4pt}}
  %\vskip0.1cm
}

% Footer estilo SimplePlus: minimalista
\setbeamertemplate{footline}{
\hfill\insertframenumber/\inserttotalframenumber\hspace{0.4cm}
}

\setbeamertemplate{itemize item}{\color{uaemcolor}\textbullet}
\setbeamertemplate{itemize subitem}{\color{uaemcolor}--}

\hypersetup{colorlinks=true, linkcolor=uaemcolor, urlcolor=uaemcolor}

%================================================
\begin{document}
%================================================

%------------------------------------------------
% PORTADA estilo SimplePlus
%------------------------------------------------
{
\setbeamertemplate{footline}{}
\setbeamertemplate{frametitle}{}
\begin{frame}


  \begin{center}
    {\small{Universidad Autónoma Del Estado De Morelos}}\\[0.05cm]
    {\footnotesize Instituto de Investigación en Ciencias Básicas y Aplicadas}
  \end{center}

\vspace{0.8cm}
  \textcolor{gray}{\rule{\textwidth}{0.4pt}}
  \vspace{-0.4cm}
  
  \begin{center}
    {\Large\bf Predicción de Calidad del Aire}\\[0.2cm]
    {\small\textit{Propuesta de Proyecto}}
  \end{center}

  \vspace{-0.4cm}
  \textcolor{gray}{\rule{\textwidth}{0.4pt}}
  \vspace{0.2cm}




    \vspace{1cm}
  \footnotesize\centering{Aprendizaje Automático -- Junio 2026\\[0.1cm]
Rebecca Castillo \\[0.1cm]
}

\end{frame}
}
%------------------------------------------------
% SLIDE: Resumen
%------------------------------------------------
\begin{frame}{Resumen de la Propuesta}
\vspace{0.2cm}
\begin{columns}[T]
\column{0.30\textwidth}
\begin{tcolorbox}[colback=uaemcolor!6, colframe=uaemcolor!60, title=\textbf{Dataset}]
\small\centering
Global Air Quality\\
\& Deforestation\\[4pt]
35{,}000 registros\\
15 variables\\
2014--2023\\[4pt]
\footnotesize Kaggle — Apache 2.0
\end{tcolorbox}

\column{0.30\textwidth}
\begin{tcolorbox}[colback=verdeenv!6, colframe=verdeenv!60, title=\textbf{Objetivos}]
\small\centering
\textbf{Opción A:}\\Predecir AQI\\(Regresión)\\[6pt]
\textbf{Opción B:}\\Clasificar categoría\\de riesgo\\(Clasificación)
\end{tcolorbox}

\column{0.30\textwidth}
\begin{tcolorbox}[colback=naranjaaqi!6, colframe=naranjaaqi!60, title=\textbf{Valor agregado}]
\small\centering
Análisis para las\\5 ciudades mexicanas\\dentro de un\\contexto global\\de 45 países
\end{tcolorbox}
\end{columns}

\end{frame}
%------------------------------------------------
% SLIDE: Dataset
%------------------------------------------------
\begin{frame}{Dataset: Global Air Quality \& Deforestation}
\begin{columns}[T]
\column{0.45\textwidth}
\vspace{-0.5cm}
\begin{tcolorbox}[colback=grisclaro, colframe=uaemcolor!60, title=\textbf{Descripción general}]
\small
\begin{tabular}{ll}
\textbf{Fuente:}    & Kaggle (Apache 2.0) \\
\textbf{Registros:} & 35{,}000 \\
\textbf{Periodo:}   & 2014 -- 2023 \\
\textbf{Ciudades:}  & 225 en 45 países \\
\textbf{Tamaño:}    & $\sim$3.1 MB (CSV) \\
\end{tabular}
\end{tcolorbox}

\vspace{0.2cm}
\begin{tcolorbox}[colback=verdeenv!6, colframe=verdeenv!60, title=\textbf{México en el dataset}]
\small
5 ciudades incluidas directamente:\\[3pt]
CDMX $\cdot$ Guadalajara $\cdot$ Monterrey\\
Puebla $\cdot$ Tijuana\\[3pt]
\footnotesize\textit{Emerging Tier}: junto a Brasil, Colombia,\\Rusia y 16 países más.
\end{tcolorbox}

\column{0.58\textwidth}
\vspace{-0.3cm}
\begin{tcolorbox}[colback=grisclaro, colframe=black!25, title=\textbf{Variables (15 columnas)}]
\small
\begin{tabular}{ll}
\vspace{0.1cm}
\textbf{Variable} & \textbf{Descripción} \\
\vspace{0.1cm}
\texttt{AQI}                  & Índice calidad del aire \\
\texttt{PM2.5 / PM10}         & Partículas suspendidas \\
\texttt{CO2\_Emissions}       & Emisiones CO\textsubscript{2} \\
\texttt{Deforestation\_Rate}  & Tasa deforestación (\%) \\
\texttt{Vehicles\_Increase}   & Crecim.\ vehicular (\%) \\
\texttt{Industries\_Increase} & Crecim.\ industrial (\%) \\
\texttt{Env\_Budget}          & Presupuesto ambiental \\
\texttt{Population\_Density}  & Densidad poblacional \\
\texttt{Green\_Space\_Ratio}  & Espacios verdes (\%) \\
\texttt{Life\_Expectancy}     & Esperanza de vida \\

\end{tabular}
\end{tcolorbox}
\end{columns}
\end{frame}

%------------------------------------------------
% SLIDE: Variable objetivo
%------------------------------------------------
\begin{frame}{Variable Objetivo: Índice de Calidad del Aire (AQI)}
\begin{columns}[T]
\column{0.44\textwidth}
\vspace{-0.5cm}
\begin{tcolorbox}[colback=grisclaro, colframe=uaemcolor!60, title=\textbf{¿Qué es el AQI?}]
\small
El \textbf{Air Quality Index} cuantifica la contaminación del aire en una escala de 0 a 500. A mayor valor, peor calidad y mayor riesgo para la salud.\\[4pt]
\footnotesize\textit{Estadísticas del dataset:}\\
Min: 15.0 \quad Max: 500.0 \quad Media: $\approx$145.0
\end{tcolorbox}

\vspace{0.2cm}
\begin{tcolorbox}[colback=grisclaro, colframe=black!25, title=\textbf{Features predictoras}]
\small
Vehículos $\cdot$ Industria $\cdot$ Densidad\\
Presupuesto ambiental $\cdot$ Deforestación\\
Espacios verdes $\cdot$ Emisiones CO\textsubscript{2}
\end{tcolorbox}

\column{0.52\textwidth}
\vspace{-0.5cm}
\begin{tcolorbox}[colback=grisclaro, colframe=black!25, title=\textbf{Categorías del AQI}]
\small
\begin{tabular}{cll}
\toprule
\textbf{Rango} & \textbf{Categoría} & \textbf{Riesgo} \\
\midrule
0--50   & {\color{verdeenv}\textbf{Bueno}}                & Mínimo \\
51--100  & {\color{yellow!60!black}\textbf{Moderado}}      & Bajo \\
101--150 & {\color{orange!80!black}\textbf{No saludable*}} & Sensibles \\
151--200 & {\color{red!70!black}\textbf{No saludable}}     & General \\
201--300 & {\color{purple!70!black}\textbf{Muy peligroso}} & Alerta \\
301--500 & {\color{black}\textbf{Peligroso}}               & Emergencia \\
\bottomrule
\end{tabular}
\vspace{3pt}
\footnotesize *Adultos mayores, niños, personas con asma.
\end{tcolorbox}

\vspace{0.15cm}
\begin{tcolorbox}[colback=uaemcolor!5, colframe=uaemcolor!40]
\footnotesize\centering
La media del dataset ($\approx$145) cae en \textbf{No saludable para grupos sensibles}, lo que da relevancia práctica al problema.
\end{tcolorbox}
\end{columns}
\end{frame}

%------------------------------------------------
% SLIDE: Dos objetivos
%------------------------------------------------
\begin{frame}{Objetivos — Dos Enfoques Propuestos}
\begin{columns}[T]
\column{0.48\textwidth}
\vspace{-0.5cm}
\begin{tcolorbox}[colback=uaemcolor!5, colframe=uaemcolor!70, title={\textbf{Opción A — Regresión}}]
\small
\textbf{Pregunta:}\\
¿Cuál será el AQI de una ciudad dado su contexto ambiental, vehicular e industrial?\\[5pt]
\textbf{Target:} \texttt{AQI} (continuo, 0--500)\\[5pt]
\textbf{Modelos:}
\begin{itemize}\small
    \item Regresión Lineal \textit{(baseline)}
    \item Random Forest Regressor
\end{itemize}
\textbf{Métricas:} RMSE, $R^2$\\[3pt]
\footnotesize\textit{Cuantifica exactamente qué tan contaminada estará una ciudad.}
\end{tcolorbox}

\column{0.48\textwidth}
\vspace{-0.5cm}
\begin{tcolorbox}[colback=verdeenv!5, colframe=verdeenv!70, title={\textbf{Opción B — Clasificación}}]
\small
\textbf{Pregunta:}\\
¿A qué categoría de riesgo pertenece una ciudad según sus indicadores ambientales?\\[5pt]
\textbf{Target:} Categoría AQI (6 clases)\\[5pt]
\textbf{Modelos:}
\begin{itemize}\small
    \item Regresión Logística \textit{(baseline)}
    \item Random Forest Classifier
\end{itemize}
\textbf{Métricas:} Accuracy, F1-score\\[3pt]
\footnotesize\textit{Útil para políticas públicas: identificar ciudades en riesgo alto.}
\end{tcolorbox}
\end{columns}

\vspace{0.2cm}
\centering\footnotesize{Ambos enfoques usan el mismo dataset y las mismas features.}

\end{frame}

%------------------------------------------------
% SLIDE: Cronograma
%------------------------------------------------
\begin{frame}{Cronograma de Trabajo}
\vspace{-0.1cm}
\begin{columns}[T]
\column{0.34\textwidth}
\begin{tcolorbox}[colback=uaemcolor!5, colframe=uaemcolor!60, title={\footnotesize\textbf{Semana 1 — Exploración}}]
\footnotesize
\begin{tabular}{ll}
\textbf{Lun 9}    & Propuesta (hoy) \\[2pt]
\textbf{Mar--Mié} & Carga de datos \\
                  & EDA inicial \\[2pt]
\textbf{Jue}      & Visualizaciones \\
                  & Correlaciones \\[2pt]
\textbf{Vie}      & Limpieza y \\
                  & preprocesamiento \\
\end{tabular}
\end{tcolorbox}

\column{0.33\textwidth}
\begin{tcolorbox}[colback=verdeenv!5, colframe=verdeenv!60, title={\footnotesize\textbf{Semana 2 — Modelado}}]
\footnotesize
\begin{tabular}{ll}
\textbf{Lun--Mar} & Entrenamiento \\
                  & de modelos \\[2pt]
\textbf{Mié}      & Evaluación y \\
                  & comparación \\[2pt]
\textbf{Jue--Vie} & Ajuste de \\
                  & hiperparámetros \\
\end{tabular}
\end{tcolorbox}

\column{0.32\textwidth}
\begin{tcolorbox}[colback=naranjaaqi!5, colframe=naranjaaqi!60, title={\footnotesize\textbf{Semana 3 — Análisis}}]
\footnotesize
\begin{tabular}{ll}
\textbf{Lun--Mar} & Análisis de \\
                  & resultados \\[2pt]
\textbf{Mié}      & Enfoque en \\
                  & ciudades MX \\[2pt]
\textbf{Jue--Vie} & Redacción del \\
                  & reporte final \\
\end{tabular}
\end{tcolorbox}
\end{columns}

\vspace{0.2cm}
\begin{tcolorbox}[colback=grisclaro, colframe=black!25, title=\footnotesize\textbf{Herramientas}]
\centering\footnotesize
\textbf{Python:} pandas, scikit-learn, matplotlib, seaborn
\quad$\cdot$\quad \textbf{Entorno:} Jupyter Notebook
\quad$\cdot$\quad \textbf{Datos:} CSV $\sim$3.1 MB
\end{tcolorbox}
\end{frame}



\end{document}
